\documentclass[12pt,a4paper]{ctexart}

% ---- 基本包 ----
\usepackage{geometry}
\geometry{left=2.5cm, right=2.5cm, top=2.5cm, bottom=2.5cm}
\usepackage{amsmath, amssymb, bm}
\usepackage{booktabs}
\usepackage{graphicx}
\usepackage{hyperref}
\usepackage{enumitem}
\usepackage{xcolor}
\usepackage{caption}
\usepackage{float}

\hypersetup{
  colorlinks=true,
  linkcolor=blue!70!black,
  citecolor=blue!70!black,
  urlcolor=blue!60!black,
}

% ---- 标题 ----
%\title{\textbf{周报：EntropyGate \\ 熵门控自适应对比解码缓解 VLM 幻觉}}
%\author{李嘉琦}
%\date{2026年3月15日 \(\sim\) 3月21日}

\begin{document}
\maketitle

\noindent\textbf{model：}LLaVA-1.5-7B-HF \quad|\quad \textbf{dataset：}CHAIR（500 samples，COCO val2014）

%======================================================================
\section{related work}
%======================================================================

VLM 的对象幻觉问题近年来受到广泛关注，对比解码是其中一类主流方案。
其基本思路是构造一个"容易产生幻觉"的辅助分布，在解码时从原始分布中减去该辅助分布的影响。
按辅助分布的构造方式，代表方法如表~\ref{tab:related}所示。

\begin{table}[H]
\centering
\caption{对比解码相关工作}
\label{tab:related}
\begin{tabular}{lll}
\toprule
\textbf{方法} & \textbf{出处} & \textbf{辅助分布构造} \\
\midrule
VCD   & CVPR 2024   & 对输入图像加高斯噪声 \\
DoLa  & ICLR 2024   & 取模型早期层的 logits \\
ACD   & EMNLP 2024  & 类似 VCD，用 $H_{\text{orig}}/(H_{\text{orig}}+H_c)$ 做自适应权重 \\
SID   & ICLR 2025   & 选择性扰动注意力中的视觉 token \\
CRoPS & 2024        & 两步嵌套：先去语言先验，再去统计偏差 \\
OPERA & CVPR 2024   & 注意力过度信任惩罚 + 回溯 \\
\bottomrule
\end{tabular}
\end{table}

这些方法普遍存在一个问题：对比强度 $\alpha$ 要么是手动设定的常数，要么依赖简单的衰减策略，没有考虑模型在生成每个 token 时的实际不确定性。

%======================================================================
\section{出发点}
%======================================================================

在跑 baseline 的过程中，我注意到模型对不同 token 的确信程度差异很大。
比如生成介词"in""on"时，输出熵很低，这种时候如果仍施加较强的对比修正，反而在本来正确的预测上引入噪声；
而在判断某个物体是否存在时，模型熵较高，固定的 $\alpha$ 又可能不够用。

让对比强度跟着当前 token 的熵走——熵高时加大力度、熵低时减小干预。

%======================================================================
\section{初始方案与 Baseline}
%======================================================================

\subsection{Baseline}

\begin{table}[H]
\centering
\caption{Baseline 结果}
\label{tab:baseline}
\begin{tabular}{lccc}
\toprule
\textbf{方法} & \textbf{CHAIRs$\downarrow$} & \textbf{CHAIRi$\downarrow$} & \textbf{Recall$\uparrow$} \\
\midrule
Vanilla                    & 52.6 & 16.3 & 72.9 \\
CRoPS ($\alpha$=1.0, $\lambda$=0.01) & 37.4 & 10.3 & 72.8 \\
\bottomrule
\end{tabular}
\end{table}

\subsection{初始 EntropyGate（Flat 结构）}

最初的设计比较直接，把熵门控以平坦的方式加在视觉和文本两路对比信号上：
\begin{align}
g_{\text{vis}} &= \sigma\!\bigl((H_t - \eta_{\text{vis}})/\tau\bigr) \cdot \alpha_{\text{base,vis}} \\
g_{\text{txt}} &= \text{decay}(t) \cdot \sigma\!\bigl((H_t - \eta_{\text{txt}})/\tau\bigr) \cdot \alpha_{\text{base,txt}} \\
\text{final} &= (1 + g_{\text{vis}} + g_{\text{txt}}) \log p - g_{\text{vis}} \log p_{\text{stat}} - g_{\text{txt}} \log p_{\text{lang}}
\end{align}

跑出来 CHAIRs=49.0, CHAIRi=12.5, Recall=77.6。Recall 不错，但幻觉抑制比 CRoPS 差了一大截。

%======================================================================
\section{实验过程与改进}
%======================================================================

\subsection{加 Floor 下限}

查看 gate 的统计数据后发现 $g_{\text{vis}}$ 的平均值只有 0.10，而 CRoPS 相当于 $\alpha$ 恒为 1.0，差了一个数量级。
于是给 $g_{\text{vis}}$ 和 $g_{\text{txt}}$ 加了下限。

最优配置 D-0.5/0.2：CHAIRs=42.6（比初始版降了 6.4），但跟 CRoPS 还有距离。

\subsection{校准 $\eta$ 阈值}

在 D 的基础上把 $\eta$ 校准到实际熵分布范围，同时把 $\alpha_{\text{base}}$ 拉大到 1.5。

最优 c1：CHAIRs=42.2, CHAIRi=10.7。CHAIRi 已经接近 CRoPS（10.3），但 CHAIRs 还差 4.8。

\subsection{根因分析}

到这一步我意识到光调超参可能不够，于是把 CRoPS 和 EntropyGate 的公式做了展开对比。CRoPS 的两步嵌套结构是：
\begin{align}
\text{intermediate} &= \frac{\log p}{\gamma} - \frac{1-\gamma}{\gamma} \log p_{\text{lang}} \\
\text{final} &= (1+\alpha) \cdot \text{intermediate} - \alpha \cdot \log p_{\text{stat}}
\end{align}

展开后原始分布前面的系数是 $(1+\alpha)/\gamma$，随时间步增长会变得很大。
而 flat 公式里原始分布的系数只有 $1 + g_{\text{vis}} + g_{\text{txt}}$，数值上只有 CRoPS 的 63\% 左右；lang prior 的对比力度更是只有 CRoPS 的 23\%。

结论很明确：\textbf{公式结构本身就有问题}，嵌套结构天然带来的交叉放大效应，flat 结构怎么调参都模拟不出来。

\subsection{嵌套 + 熵门控（主要结果）}

想清楚之后做法就很自然了：复用 CRoPS 的嵌套结构，只用熵门控替换第二步里固定的 $\alpha$：
\begin{align}
\text{intermediate} &= \log p + \frac{1-\gamma_t}{\gamma_t}\bigl(\log p - \log p_{\text{lang}}\bigr) \label{eq:nested_step1} \\
g_{\text{vis}} &= \alpha_{\min} + (\alpha_{\text{base}} - \alpha_{\min}) \cdot \sigma\!\left(\frac{H_t - \eta}{\tau}\right) \label{eq:entropy_gate} \\
\text{final} &= (1 + g_{\text{vis}}) \cdot \text{intermediate} - g_{\text{vis}} \cdot \log p_{\text{stat}} \label{eq:nested_step2}
\end{align}

\begin{table}[H]
\centering
\caption{Scheme E 结果}
\label{tab:scheme_e}
\begin{tabular}{lcccccc}
\toprule
\textbf{配置} & $\alpha_{\text{base}}$ & $\alpha_{\min}$ & \textbf{CHAIRs$\downarrow$} & \textbf{CHAIRi$\downarrow$} & \textbf{Recall$\uparrow$} \\
\midrule
E4（固定 $\alpha$=1.0，$\approx$CRoPS） & 1.0 & 1.0 & 37.0 & 10.1 & 73.2 \\
\textbf{E5（熵门控 [0.5, 1.5]）} & \textbf{1.5} & \textbf{0.5} & \textbf{35.6} & \textbf{9.9} & \textbf{73.6} \\
\bottomrule
\end{tabular}
\end{table}

E4 把 $\alpha$ 固定为 1.0，等价于 CRoPS，跑出 37.0，和 CRoPS 的 37.4 基本一致，说明嵌套结构对接没有问题。
E5 加上熵门控后 CHAIRs 降到 35.6，这 1.4 的提升就是熵门控的纯贡献。


\subsection{Latent Space 方向}

把对比操作搬到隐状态空间。
主要考虑是：模型最后一层到 logits 之间有个 RMSNorm，它是非线性的，所以在 hidden state 空间做对比再投影回来，和直接在 logit 空间做对比得到的结果是不一样的。
实现了三种方法：

\paragraph{HSC（Hidden State Contrastive）。}
将 stat bias 对比搬到 hidden state 空间，lang prior 对比留在 logit 空间：
\begin{align}
\mathbf{h}_{\text{contrasted}} &= \mathbf{h}_{\text{orig}} + g_{\text{vis}}(H_t) \cdot (\mathbf{h}_{\text{orig}} - \mathbf{h}_{\text{stat}}) \\
\text{logits}_{\text{corrected}} &= W_{\text{head}} \cdot \text{RMSNorm}(\mathbf{h}_{\text{contrasted}})
\end{align}
然后在 logit 空间做 lang prior 对比（带时间衰减）。

\paragraph{LEG（Latent Entropy Gate）。}
对比公式跟 E5 一样在 logit 空间做，但驱动门控的熵信号改成从模型中间层（如第 16 层）取。中间层的不确定性可能比输出层更早反映幻觉风险。

\paragraph{LLH（Latent-Logit Hybrid）。}
两阶段都用熵门控——第一阶段在 hidden state 空间做 stat bias 对比（用原始熵 $H_t$ 门控），第二阶段在 logit 空间做 lang prior 对比（用修正后分布的熵 $H_{\text{corrected}}$ 门控）。

三种方法的代码和实验脚本均已完成，各配了 5 组参数（包括参考 LASER 的高 $\eta$ 配置）。

\subsubsection{HSC 结果}

\begin{table}[H]
\centering
\caption{HSC（Hidden State Contrastive）结果}
\label{tab:hsc}
\begin{tabular}{lccccccc}
\toprule
\textbf{配置} & $\alpha_{\text{base}}$ & $\alpha_{\min}$ & $\eta$ & $\tau$ & \textbf{CHAIRs$\downarrow$} & \textbf{CHAIRi$\downarrow$} & \textbf{Recall$\uparrow$} \\
\midrule
H1 & 1.0 & 0.3 & 0.10 & 0.05 & \textbf{43.8} & 13.4 & 73.4 \\
H2 & 1.5 & 0.3 & 0.10 & 0.05 & 44.4 & 13.6 & 73.9 \\
H3 & 1.5 & 0.5 & 0.10 & 0.05 & 44.6 & \textbf{13.3} & 72.8 \\
H4 & 1.5 & 0.3 & 0.30 & 0.10 & 46.2 & 12.9 & \textbf{74.8} \\
H5 & 1.0 & 0.3 & 0.50 & 0.15 & 48.0 & 13.2 & 74.1 \\
\bottomrule
\end{tabular}
\end{table}

HSC 最优 H1：CHAIRs=43.8。整体不如 logit 空间方案，RMSNorm 的非线性未带来预期收益。

\subsubsection{LEG 结果}

\begin{table}[H]
\centering
\caption{LEG（Latent Entropy Gate）结果}
\label{tab:leg}
\begin{tabular}{lcccccccc}
\toprule
\textbf{配置} & \textbf{Layer} & $\alpha_{\text{base}}$ & $\alpha_{\min}$ & $\eta$ & $\tau$ & \textbf{CHAIRs$\downarrow$} & \textbf{CHAIRi$\downarrow$} & \textbf{Recall$\uparrow$} \\
\midrule
L1 & $-$16 & 1.5 & 0.5 & 0.10 & 0.05 & 39.2 & 10.8 & 72.3 \\
L2 & $-$24 & 1.5 & 0.5 & 0.10 & 0.05 & 39.0 & 10.7 & 72.3 \\
L3 & $-$8  & 1.5 & 0.5 & 0.10 & 0.05 & 39.0 & \textbf{10.3} & \textbf{73.4} \\
L4 & $-$16 & 2.0 & 0.3 & 0.10 & 0.05 & 39.0 & 10.5 & 72.5 \\
\textbf{L5} & $-$16 & 1.5 & 0.3 & 0.30 & 0.10 & \textbf{38.6} & 10.6 & 72.0 \\
\bottomrule
\end{tabular}
\end{table}

LEG 最优 L5：CHAIRs=38.6。中间层熵信号有效，各层差异不大，LASER 风格高 $\eta$ 略优。但仍不如 E5（35.6）。

\subsubsection{LLH 结果}

\begin{table}[H]
\centering
\caption{LLH（Latent-Logit Hybrid）结果}
\label{tab:llh}
\begin{tabular}{lccccccc}
\toprule
\textbf{配置} & $\alpha_{\text{base}}$ & $\alpha_{\min}$ & $\eta$ & $\tau$ & \textbf{CHAIRs$\downarrow$} & \textbf{CHAIRi$\downarrow$} & \textbf{Recall$\uparrow$} \\
\midrule
\textbf{J1} & 1.0 & 0.3 & 0.10 & 0.05 & \textbf{45.8} & 13.5 & 74.3 \\
J2 & 1.5 & 0.3 & 0.10 & 0.05 & 45.6 & 12.9 & 74.3 \\
J3 & 1.5 & 0.5 & 0.10 & 0.05 & 45.6 & \textbf{12.8} & 73.5 \\
J4 & 1.5 & 0.3 & 0.30 & 0.10 & 49.6 & 14.9 & 74.5 \\
J5 & 1.0 & 0.3 & 0.50 & 0.15 & 50.6 & 14.6 & \textbf{75.9} \\
\bottomrule
\end{tabular}
\end{table}

LLH 最优 J3：CHAIRs=45.6。双阶段熵门控反而不如单阶段，hidden state 对比引入的非线性在此场景下是负面的。

%======================================================================
\section{当前最优结果}
%======================================================================

\subsection{总览对比}

\begin{table}[H]
\centering
\caption{总览对比}
\label{tab:final}
\begin{tabular}{lcccc}
\toprule
\textbf{方法} & \textbf{CHAIRs$\downarrow$} & \textbf{CHAIRi$\downarrow$} & \textbf{Recall$\uparrow$} & \textbf{vs CRoPS} \\
\midrule
Vanilla & 52.6 & 16.3 & 72.9 & — \\
CRoPS   & 37.4 & 10.3 & 72.8 & baseline \\
\textbf{EntropyGate E5} & \textbf{35.6} & \textbf{9.9} & \textbf{73.6} & \textbf{$-$1.8 / $-$0.4 / $+$0.8} \\
\bottomrule
\end{tabular}
\end{table}

三项指标全面优于 CRoPS。

\subsection{全部方案汇总表}

\begin{table}[H]
\centering
\caption{全部实验方案汇总（4bit 量化，CHAIR 500 samples）}
\label{tab:all_methods}
\small
\begin{tabular}{llcccc}
\toprule
\textbf{类别} & \textbf{方法/配置} & \textbf{CHAIRs$\downarrow$} & \textbf{CHAIRi$\downarrow$} & \textbf{Recall$\uparrow$} & \textbf{公式结构} \\
\midrule
\multicolumn{6}{l}{\textit{Baseline}} \\
& Vanilla (Greedy) & 52.6 & 16.3 & 72.9 & — \\
& CRoPS ($\alpha$=1.0) & 37.4 & 10.3 & 72.8 & nested, 固定 $\alpha$ \\
\midrule
\multicolumn{6}{l}{\textit{Flat 结构（Logit 空间）}} \\
& 原始 EntropyGate & 49.0 & 12.5 & \textbf{77.6} & flat, 熵门控 \\
& Scheme D best (0.5/0.2) & 42.6 & 11.0 & 76.5 & flat + floor \\
& Scheme D+A best (c1) & 42.2 & 10.7 & 75.3 & flat + floor + adaptive \\
& Dir12 best (combo) & 40.2 & 10.9 & 73.3 & flat, 综合调优 \\
& Dir34 best (abtxt2) & 41.8 & 11.2 & 74.7 & flat + adaptive\_eta \\
\midrule
\multicolumn{6}{l}{\textit{Nested 结构（Logit 空间）}} \\
& Scheme E4 (固定 $\alpha$=1.0) & 37.0 & 10.1 & 73.2 & nested, $\approx$CRoPS \\
& \textbf{Scheme E5 (熵门控 [0.5,1.5])} & \textbf{35.6} & \textbf{9.9} & 73.6 & nested + 熵门控 \\
& Scheme F best (ACD 熵比值) & 39.2 & 10.0 & 76.0 & nested + 熵比值 \\
& Scheme G5 (E5 + 对齐参数) & 35.6 & 9.9 & 73.6 & nested + 对齐参数 \\
\midrule
\multicolumn{6}{l}{\textit{Latent Space 方向}} \\
& HSC best (H1) & 43.8 & 13.4 & 73.4 & hidden state 对比 \\
& LEG best (L5) & 38.6 & 10.6 & 72.0 & 中间层熵门控 \\
& LLH best (J3) & 45.6 & 12.8 & 73.5 & hidden + logit 混合 \\
\midrule
\multicolumn{6}{l}{\textit{VCD 方向（do\_sample，不可直接对比）}} \\
& VCD & 58.4 & 17.4 & 75.6 & VCD \\
& VCD + EntropyGate & 55.8 & 17.4 & 75.7 & VCD + 熵门控 \\
\bottomrule
\end{tabular}
\end{table}

核心结论：
\begin{enumerate}[leftmargin=2em]
\item \textbf{嵌套结构是决定性因素}：flat$\to$nested 带来 $\sim$5 个点的 CHAIRs 提升，远超任何超参数调优（flat 天花板 $\approx$40，nested E5=35.6）。
\item \textbf{熵门控在 nested 上仍有贡献}：E4(固定$\alpha$=1.0)=37.0 $\to$ E5(熵门控)=35.6，纯贡献 1.4。
\item \textbf{E 和 G 完全一致}：cutoff/safety 参数不影响结果。
\item \textbf{ACD 熵比值（F）保护 Recall 但抑制力度不足}：CHAIRs=39.2 vs E5=35.6。
\item \textbf{Latent Space 方向}：LEG（中间层熵信号）有一定潜力（38.6），但 HSC 和 LLH 的 hidden state 对比效果较差（43.8 / 45.6），RMSNorm 非线性在此场景下未带来正面收益。
\item \textbf{E5 三项指标全面超过 CRoPS}：CHAIRs $-$1.8, CHAIRi $-$0.4, Recall $+$0.8。
\end{enumerate}

%======================================================================
\section{新论文}
%======================================================================

\subsection{LEAD（CVPR 2026）}

论文全称 \textit{Thinking in Uncertainty: Mitigating Hallucinations in MLRMs with Latent Entropy-Aware Decoding}（arXiv: 2603.13366）。

LEAD 观察到转折词（because, however 等）常处于高熵状态，幻觉容易在这之后出现。
它的做法有两点：一是高熵时不选具体 token，改用所有 token embedding 的概率加权混合作为下一步输入：
\begin{equation}
\tilde{\mathbf{e}}_t = \begin{cases}
\mathbf{e}(r_t), & H_t < \hat{H} \\
\mathbb{E}_{v \sim p_t}[\mathbf{e}(v)], & H_t \geq \hat{H}
\end{cases}
\end{equation}

二是高熵时注入视觉锚点：
\begin{equation}
\tilde{\mathbf{e}}_t^* = (1-\lambda)\,\mathbb{E}_{v \sim p_t^*}[\mathbf{e}(v)] + \lambda\,\mathbf{e}_{\text{vis}}, \quad \lambda = 0.4
\end{equation}

在 R1-Onevision-7B 上，MMHalu $+$4.7\%, VStar $+$4.7\%, MathVista $+$2.3\%。

\subsection{LASER}

LASER 认为 CoT 的显式文本推理存在信息瓶颈。
它的做法是在隐状态空间保持"语义叠加态"，用动态窗口上的 soft distribution 做监督，并在模型不确定时注入 hard target：
\begin{equation}
P_t^{\text{target}} = \begin{cases}
\alpha \cdot y_{\text{hard}} + (1-\alpha) \cdot Q_t, & H(Q_t) > \eta \\
Q_t, & H(Q_t) \leq \eta
\end{cases}
\end{equation}
其中 $\eta=0.6$，约 10\% 的 token 触发干预。HallusionBench 提升 11.36\%，推理 token 量减少 97.3\%。


\begin{table}[H]
\centering
\caption{CRoPS vs LEAD 对比}
\label{tab:lead_vs_crops}
\begin{tabular}{lll}
\toprule
 & \textbf{CRoPS} & \textbf{LEAD} \\
\midrule
干预方式 & 修改 logits 分布 & 切换输入 embedding \\
干预时机 & 每个 token & 仅高熵 token \\
目标场景 & 通用 VLM & 推理模型（R1 系列） \\
核心假设 & 幻觉来自视觉/语言偏差 & 幻觉来自高熵点错误坍缩 \\
\bottomrule
\end{tabular}
\end{table}


%======================================================================
\section{反思}
%======================================================================

\subsection{反思}

\begin{enumerate}[leftmargin=2em]
\item 前期在 flat 结构上花了不少时间调超参，回头看应该更早去分析公式结构本身。flat 怎么调都过不了 40，换成 nested 一下子到了 35.6。
\item 一个意外的现象：fp16 下 nested CHAIRs=55.4，4-bit 量化下是 35.6，差了将近 20 个点。猜测是量化引入的噪声反而让对比信号更有效。
\item "用熵门控替换固定 $\alpha$"这个想法不局限于 CRoPS，理论上 VCD、DoLa、SID 都能用，可以定位为对比解码的通用自适应层。
\item 从 logit 到 latent 的推进比较自然——RMSNorm 的非线性让 hidden state 空间的对比有独立意义，中间层的熵信号也可能比输出层更早反映幻觉风险。
\end{enumerate}

\end{document}
